\section{Four aggregations for three strict inequalities under PDLC}
\label{sec:four-aggregation}

The preceding example needs infinitely many good aggregations and does
not satisfy PDLC for its homogeneous matrices. We now establish the
sharp finite bound for three strict inequalities under that condition.
Here PDLC means that $Q_\theta\succ0$ for some $\theta\in\R^3$;
the coefficients of $\theta$ may have either sign.

\begin{theorem}[Four aggregations for strict PDLC systems]
\label{thm:strict-four}
Let $n\geq1$ and $m=3$. Suppose that $S\ne\varnothing$,
$\conv S\ne\R^n$, and $Q_1,Q_2,Q_3$ satisfy PDLC. Then there
are $1\leq\ell\leq4$ good multipliers $\lambda^1,\ldots,\lambda^\ell$
such that
\begin{equation}\label{eq:strict-four}
 \conv S=\bigcap_{j=1}^{\ell}\{x:f_{\lambda^j}(x)<0\}.
\end{equation}
Each $Q_{\lambda^j}$ has exactly one negative eigenvalue.
The original matrices need not be linearly independent, and the original
nonstrict feasible set need not be bounded or the closure of its interior.
For every $n\geq3$, four is necessary in some systems.
\end{theorem}

Blekherman and Dunbar's four-bound for regular nonstrict systems
\citep[Theorem~1.4]{BD2025,BDpreprint} is the substantive external
input. We transfer that bound to arbitrary strict systems by compact
inward perturbation and a limit that preserves strict validity.
Blekherman, Dey, and Sun previously gave an upper bound of six
\citep[Corollary~2.20]{BDSpreprint}; their Example~2.21 supplies the
sharpness construction below. Thus Theorem~\ref{thm:strict-four}
rules out the six-necessary example proposed in their Conjecture~3.2.
The number four and its lower-bound example are established antecedents;
the extension here concerns the unrestricted strict setting.

Dunbar's dissertation also states a four-bound for regular nonstrict
sets without the no-points-at-infinity hypothesis
\citep[Theorem~5.0.5]{Dunbar2025Thesis}. We acknowledge that stronger
prior statement and use only the fully qualified published theorem below.
To the best of our knowledge, the complete transfer to arbitrary strict
systems in Theorem~\ref{thm:strict-four} has not previously been given.
This claim concerns that transfer, including possible dependence of the
original matrices, rather than the first statement of a four-bound without
an infinity assumption. The negative-eigenvector limit used below has an
antecedent for a fixed feasible set in the proof of
\citet[Proposition~3.15, Section~8]{BDpreprint}.

\subsection{Regular inward perturbations}

We state the external input with its standing independence assumption.
For a triple $\widetilde Q_i$ with leading blocks $\widetilde A_i$,
the phrase \emph{no points at infinity} means
\begin{equation}\label{eq:no-infinity}
 \{d\in\R^n:d\transpose\widetilde A_i d\leq0\text{ for all }i\}
 =\{0\}.
\end{equation}

\begin{theorem}[Blekherman--Dunbar, external input]
\label{thm:bd-four-input}
Let $\widetilde Q_1,\widetilde Q_2,\widetilde Q_3\in\Sym^{n+1}$
be linearly independent and satisfy PDLC, and set
$C=\{x:(x,1)\transpose\widetilde Q_i(x,1)\leq0\text{ for all }i\}$.
If $\interior C\ne\varnothing$, $C=\cl(\interior C)$, and
\eqref{eq:no-infinity} holds, then
\[
 \cl(\conv C)=\bigcap_{j=1}^{\ell}
 \left\{x:(x,1)\transpose\Bigl(\sum_i\mu_i^j\widetilde Q_i\Bigr)(x,1)
 \leq0\right\}
 \qquad(\ell\leq4),
\]
where each $\mu^j\geq0$ is nonzero and each aggregated matrix has
at most one negative eigenvalue.
\end{theorem}

This is \citet[Theorem~1.4]{BD2025,BDpreprint}; no smoothness
assumption on the spectral curve is needed. The proof in the cited
preprint explicitly includes $n=1,2$ in Proposition~8.7 and combines
Propositions~8.6, 8.7, and 8.10. Consequently the input applies in
every positive dimension used here.

\begin{lemma}[Regular sublevels]\label{lem:regular-levels}
Let $h:M\to\R$ be continuous on a space with a countable base.
For all but at most countably many $c\in\R$,
\[
 \{h\leq c\}=\cl_M\{h<c\}.
\]
\end{lemma}
\begin{proof}
Continuity gives the inclusion from right to left. If it is strict,
there is $z$ with $h(z)=c$ and an open neighborhood $U$ of $z$
containing no point where $h<c$. Choose a basic open set $B$ with
$z\in B\subseteq U$. Then $c=\inf_{y\in B}h(y)$. There are at
most countably many such real infima, one for each basic open set.
\end{proof}

\begin{lemma}[Compact inward systems]\label{lem:inward-pdlc}
Under the hypotheses of Theorem~\ref{thm:strict-four}, there are
positive definite matrices $R_1,R_2,R_3$ and a decreasing sequence
$\varepsilon_k\downarrow0$ such that, on defining
\begin{align*}
 Q_i^\varepsilon&=Q_i+\varepsilon R_i,
 &w_i(x)&=(x,1)\transpose R_i(x,1),\\
 S_\varepsilon&=\{x:f_i(x)+\varepsilon w_i(x)<0\ \forall i\},
 &C_\varepsilon&=\{x:f_i(x)+\varepsilon w_i(x)\leq0\ \forall i\},
\end{align*}
each selected triple satisfies every hypothesis of
Theorem~\ref{thm:bd-four-input}. Moreover, $C_{\varepsilon_k}$ is
compact, $C_{\varepsilon_k}=\cl S_{\varepsilon_k}$,
$S_{\varepsilon_k}\subseteq S$, and every fixed finite subset of $S$
is contained in $S_{\varepsilon_k}$ for all sufficiently large $k$.
\end{lemma}
\begin{proof}
Since $n+1\geq2$, take, with standard basis vectors $e_1,e_2$,
\[
 R_1=I,\qquad R_2=I+e_1e_1\transpose,\qquad
 R_3=I+(e_1+e_2)(e_1+e_2)\transpose.
\]
These are positive definite and linearly independent: in a vanishing
linear combination the $(1,2)$ entry first forces the coefficient of
$R_3$ to vanish, and the two diagonal entries then force the others to
vanish. Write $B_i\succ0$ for their leading $n\times n$ blocks.
Lemma~\ref{lem:negative-direction} excludes a nonzero $d$ satisfying
$d\transpose A_i d<0$ for every $i$. Therefore, for every
$\varepsilon>0$, the inequalities
$d\transpose(A_i+\varepsilon B_i)d\leq0$ for all $i$ force $d=0$.
Thus the perturbed system has no points at infinity.

It is also bounded. Otherwise there would be $x_j\in C_\varepsilon$
with $\|x_j\|\to\infty$. After taking a subsequence,
$x_j/\|x_j\|\to d$ with $\|d\|=1$. Divide each perturbed
inequality by $\|x_j\|^2$ and take limits to obtain the forbidden
nonpositive leading inequalities at $d$. Closedness now makes
$C_\varepsilon$ compact.

Fix $x^0\in S$ and a PDLC certificate $\theta$. For every sufficiently
small $\varepsilon>0$, $x^0\in S_\varepsilon$ and
$Q_\theta+\varepsilon\sum_i\theta_iR_i\succ0$, by strictness and
continuity. The perturbed matrices are independent except at finitely
many parameters. Indeed, choose three coordinates of $\Sym^{n+1}$
on which the coordinate matrix of the $R_i$ has nonzero determinant.
The corresponding determinant for the $Q_i+\varepsilon R_i$ is a
cubic polynomial whose leading coefficient is that nonzero determinant.
Its zeros form a finite set. This argument does not assume independence
of the original $Q_i$.

Because every $w_i$ is positive, the continuous function
\[
 h(x)=\max_{1\leq i\leq3}\frac{f_i(x)}{w_i(x)}
\]
satisfies $S_\varepsilon=\{h<-\varepsilon\}$ and
$C_\varepsilon=\{h\leq-\varepsilon\}$. Apply
Lemma~\ref{lem:regular-levels} and choose $\varepsilon_k\downarrow0$
so that $-\varepsilon_k$ avoids its countable exceptional levels and
$\varepsilon_k$ avoids the finite dependence set, with every parameter
small enough for strict feasibility and PDLC. Such a sequence exists
because every nonempty real interval is uncountable. Then
\[
 C_{\varepsilon_k}=\cl S_{\varepsilon_k}
   =\cl(\interior C_{\varepsilon_k}),\qquad
 \interior C_{\varepsilon_k}\ne\varnothing.
\]
The second equality follows from
$S_{\varepsilon_k}\subseteq\interior C_{\varepsilon_k}
\subseteq C_{\varepsilon_k}$ and closedness. Positivity of $w_i$
gives $S_\varepsilon\subseteq S$. Finally, finitely many strict
inequalities at finitely many points remain strict for all sufficiently
small $\varepsilon$, proving the last assertion.
\end{proof}

\subsection{Strictification and the multiplier limit}

Two details are essential: a redundant nonstrict quadratic can remove
interior points when made strict, and a limiting quadratic can lose strict
validity unless its negative component is controlled.

\begin{lemma}[Strictification]\label{lem:strictification}
Let $U\subset\R^n$ be nonempty and open with bounded hull. Suppose
\[
 \cl(\conv U)=\bigcap_{j=1}^q\{x:p_j(x)\leq0\},
\]
where the $p_j$ are quadratic polynomials. After deleting every $p_j$
that is globally nonpositive, at least one inequality remains and
\[
 \conv U=\bigcap_{j\text{ retained}}\{x:p_j(x)<0\}.
\]
\end{lemma}
\begin{proof}
Deleting such polynomials preserves the nonstrict intersection. Some
inequality remains because a nonempty bounded hull is not $\R^n$.
If a quadratic $p$ vanishes at an interior point $x$ of $\{p\leq0\}$,
then $x$ is a local maximum with value zero. Its gradient vanishes and
its Hessian is negative semidefinite. The exact quadratic Taylor formula
therefore gives $p\leq0$ everywhere. For every retained polynomial this
is impossible, so $\interior\{p_j\leq0\}=\{p_j<0\}$.
The interior of a finite intersection equals the intersection of its
interiors. Finally, the nonempty open convex set $\conv U$ equals the
interior of its closure, which proves the formula.
\end{proof}

For completeness, the relevant negative components have a direct conic
description. If a symmetric matrix $M$ has exactly one negative eigenvalue
$-a<0$, choose its unit eigenvector $v$ and put $P=M+avv\transpose$.
Then $P\succeq0$, $Pv=0$, and
\begin{equation}\label{eq:negative-components}
 \{z:z\transpose Mz<0\}
 =\{z:\|P^{1/2}z\|<\sqrt a\,v\transpose z\}
  \;\cup\;
  \{z:\|P^{1/2}z\|<-\sqrt a\,v\transpose z\}.
\end{equation}
Both sets are convex, disjoint, and nonempty. This formula also covers
singular $M$ and the case $P=0$. A connected subset of the negative
sublevel lies in one of the two sets.

\begin{proof}[Proof of Theorem~\ref{thm:strict-four}, upper bound]
Use the sequence from Lemma~\ref{lem:inward-pdlc} and abbreviate
$S_k=S_{\varepsilon_k}$, $C_k=C_{\varepsilon_k}$.
Because $C_k=\cl S_k$,
\[
 \cl(\conv C_k)=\cl(\conv S_k).
\]
Indeed, each closed convex set containing $S_k$ also contains $C_k$,
and the reverse inclusion follows from $S_k\subseteq C_k$.
Theorem~\ref{thm:bd-four-input} followed by
Lemma~\ref{lem:strictification} gives a description of $\conv S_k$
by between one and four strict aggregations of the perturbed inequalities.
Every retained aggregation is good for $S_k$ and has exactly one
negative homogeneous eigenvalue, since it is strictly negative at $x^0$.
Duplicate multipliers to obtain four indices and normalize each in
the simplex $\Delta_3=\{\lambda\geq0:\sum_i\lambda_i=1\}$.
Write them as $\lambda^{k,j}$, $1\leq j\leq4$. By compactness,
after a common subsequence
$\lambda^{k,j}\to\lambda^j\in\Delta_3$ for every $j$.
Their matrices satisfy
\begin{equation}\label{eq:four-matrix-limit}
 M_{k,j}=Q_{\lambda^{k,j}}
       +\varepsilon_k\sum_i\lambda_i^{k,j}R_i
 \ \longrightarrow\ M_j=Q_{\lambda^j}.
\end{equation}
The perturbation term tends to zero because its norm is at most
$\varepsilon_k\max_i\|R_i\|$. Continuity of ordered eigenvalues
shows that $M_j$ has at most one negative eigenvalue. It has at least
one, since $f_{\lambda^j}(x^0)<0$. In particular every limiting
multiplier is nonzero and every limiting matrix has exactly one negative
eigenvalue.

Choose a unit negative eigenvector $v_{k,j}$ of $M_{k,j}$, oriented
so that $v_{k,j}\transpose(x^0,1)>0$. This scalar product cannot
vanish because the quadratic value at $(x^0,1)$ is negative. Take a
further common subsequence so that $v_{k,j}\to v_j$ for all $j$.
The unique negative eigenvalues converge to the negative eigenvalue of
$M_j$; passing to the eigenvector equation shows that $v_j$ is its
unit eigenvector.

Fix any $x\in S$. For all sufficiently large $k$, both $x$ and $x^0$
belong to $S_k$. The segment joining their lifts is strictly negative
for $M_{k,j}$, because the $k$th aggregation is good on $\conv S_k$.
The two scalar products with $v_{k,j}$ therefore have the same sign:
otherwise that segment would meet $v_{k,j}^{\perp}$, where $M_{k,j}$
is positive semidefinite. Hence $v_{k,j}\transpose(x,1)>0$.
Taking limits gives $v_j\transpose(x,1)\geq0$. Equality is
impossible, since $f_{\lambda^j}(x)<0$ and $M_j$ is positive
semidefinite on $v_j^\perp$. Thus every lift of every point of $S$
lies in the same convex negative component of $M_j$ in
\eqref{eq:negative-components}. All finite convex combinations of these
lifts remain in that component. This proves
\[
 \conv S\subseteq\bigcap_{j=1}^4\{x:f_{\lambda^j}(x)<0\}
\]
and proves strict goodness of every limit.

Conversely, a point strictly satisfying all four limiting inequalities
strictly satisfies all four perturbed inequalities for sufficiently
large $k$, by \eqref{eq:four-matrix-limit} and finiteness of the
index set. It then belongs to $\conv S_k\subseteq\conv S$.
This proves \eqref{eq:strict-four}; remove repeated multipliers if
desired. No closure is taken in this last argument.
\end{proof}

\begin{remark}[Dependent triples]\label{rem:dependent-four}
The preceding proof already covers dependence. A stronger bound of two
follows from the classical two-constraint result and the cone reduction
also noted in \citet[Section~2.5, after Example~2.21]{BDSpreprint}.
Indeed, evaluation at
a fixed feasible lift is strictly negative on every nonzero nonnegative
combination of the $Q_i$. Their finitely generated cone is therefore
pointed. If their span has dimension at most two, intersecting this cone
with the affine hyperplane on which that evaluation is $-1$ gives a
compact segment or a point. Its at most two endpoints lie on rays
generated by original matrices. Every original inequality is a nonzero
nonnegative combination of those at most two inequalities, so they define
the same $S$. Applying \citet[Theorem~1]{Yildiran2009}, with its
opposite sign convention, gives at most two good aggregations; translating
their multipliers back to the original rows preserves goodness. A single
remaining inequality may be duplicated when applying that theorem.
This observation is not needed for the four-bound.
\end{remark}

\subsection{Sharpness and the closed hull}

\begin{example}[The four-necessary example of Blekherman--Dey--Sun]
\label{ex:four-sharp}
Let $n\geq3$, $s=x_1$, and $\rho=\sum_{i=2}^n x_i^2$. Consider
\begin{equation}\label{eq:four-sharp-system}
 f_1=-s^2+1+\rho,\qquad
 f_2=s^2+5s-4+\rho,\qquad f_3=-s-\rho.
\end{equation}
This is \citet[Example~2.21]{BDSpreprint}. The identity
$-7f_1-3f_2-15f_3=4s^2+5\rho+5$ proves PDLC.
The point $(s,\rho)=(-3,4)$ is strictly feasible. Every feasible point
has $s<-1$: $f_1<0$ forces $|s|>1$, while $f_2<0$ excludes $s>1$.
Consequently the hull is proper, and the $f_1$ inequality is good because
these points lie in its negative $s$ component. The convex ball inequality
$f_2<0$ is good as well. The two sums are
\[
 f_1+f_3=-s^2-s+1,
 \qquad f_2+f_3=s^2+4s-4.
\]
The first has one negative homogeneous eigenvalue and is negative on
the interval $s<\alpha=(-1-\sqrt5)/2$ containing every feasible $s$.
The second is convex and also has one negative homogeneous eigenvalue.
Thus the following four rays are good:
\[
 r_1=(1,0,0),\quad r_2=(0,1,0),\quad
 r_3=(1,0,1),\quad r_4=(0,1,1).
\]

We verify directly that each ray is indispensable. If $(a,b,c)$ is good,
then $c\leq a+b$: otherwise the coefficient $a+b-c$ of $\rho$ gives
at least $n-1\geq2$ negative eigenvalues. Hence every good multiplier
belongs to
\begin{equation}\label{eq:four-ray-cone}
 \{(a,b,c)\geq0:c\leq a+b\}=\cone\{r_1,r_2,r_3,r_4\}.
\end{equation}
For the equality, choose $t\in[\max(0,c-b),\min(a,c)]$ and write
$(a,b,c)=(a-t)r_1+(b-c+t)r_2+tr_3+(c-t)r_4$.
Put $\beta=-2-2\sqrt2$. The four witnesses below have the displayed
slacks in the order $(f_{r_1},f_{r_2},f_{r_3},f_{r_4})$:
\[
\begin{array}{c|c|c}
 \text{designated ray}&(s,\rho)&\text{slacks}\\\hline
 r_1&(-2,3)&(0,-7,-1,-8)\\
 r_2&(-4,8)&(-7,0,-11,-4)\\
 r_3&(\alpha,0)&(\alpha,4\alpha-3,0,3\alpha-3)\\
 r_4&(\beta,0)&(4\beta-3,\beta,3\beta-3,0).
\end{array}
\]
Every off-designated slack is strictly negative, since $\alpha,\beta<0$.
Each witness is realizable by setting $x_2=\sqrt\rho$ and the other
coordinates to zero. At the witness for $r_j$, any nonzero multiplier in
\eqref{eq:four-ray-cone} is strictly negative unless it is a positive
multiple of $r_j$. But that witness is outside $\conv S$, since its
designated good inequality is zero. Every exact strict-good description
must therefore contain all four rays. The upper bound already proved
then implies that these four inequalities describe the hull exactly.
This proves the sharpness assertion in Theorem~\ref{thm:strict-four}.
The argument makes no sharpness claim for $n=1,2$.
\end{example}

The strict result also provides a closed conic description, provided the
negative components retain their orientation. The conic representation
principle is already recorded in \citet[Remark~2.19]{BDSpreprint};
the bound of four follows
from Theorem~\ref{thm:strict-four}.

\begin{corollary}[Four oriented second-order cone constraints]
\label{cor:four-soc}
Under the hypotheses of Theorem~\ref{thm:strict-four}, choose its
multipliers and write $Q_{\lambda^j}=P_j-v_jv_j\transpose$ with
$P_j\succeq0$, $P_jv_j=0$, and
$v_j\transpose(x^0,1)>0$ for one $x^0\in S$. Then
\begin{align}
 \conv S&=\{x:\|P_j^{1/2}(x,1)\|<v_j\transpose(x,1)
                    \text{ for }1\leq j\leq\ell\},\label{eq:four-soc-open}\\
 \cl(\conv S)&=\{x:\|P_j^{1/2}(x,1)\|\leq v_j\transpose(x,1)
                    \text{ for }1\leq j\leq\ell\}.\label{eq:four-soc-closed}
\end{align}
\end{corollary}
\begin{proof}
Use the spectral decomposition in \eqref{eq:negative-components},
absorbing the square root of the negative eigenvalue's magnitude into
$v_j$. Strict goodness and connectedness place $\conv S$ in the
component containing $(x^0,1)$. The intersection of those components
contains the hull and is contained in \eqref{eq:strict-four}, proving
\eqref{eq:four-soc-open}. Its closed version contains the closure.
Conversely, let $x$ satisfy every closed conic inequality and put
$x_t=(1-t)x+tx^0$, $0<t\leq1$. Convexity of the norm gives
\[
 v_j\transpose(x_t,1)-\|P_j^{1/2}(x_t,1)\|
 \geq t\bigl(v_j\transpose(x^0,1)-\|P_j^{1/2}(x^0,1)\|\bigr)>0.
\]
Thus $x_t\in\conv S$ and $x_t\to x$ as $t\downarrow0$.
\end{proof}

Replacing $<$ by $\leq$ in \eqref{eq:strict-four} without keeping
the orientations need not give this closure. Nor must this closure
equal the hull of the original nonstrict set. For example, for $n\geq2$,
take the half-ball example of \citet[Example~2.23]{BDSpreprint} and
append one redundant inequality:
\[
 f_1=-x_1^2+x_1,\qquad f_2=\|x\|^2-1,\qquad
 f_3=-\|x\|^2-1.
\]
Here $-Q_3=I$, so PDLC holds, and
$S=\conv S=\{x:x_1<0,\ \|x\|<1\}$. The first two strict
inequalities are good and describe $S$. Their nonstrict counterparts
describe the closed half-ball together with the isolated point $e_1$.
The original nonstrict set contains that same extra point. Its hull is
therefore larger than $\cl(\conv S)$.

Theorem~\ref{thm:strict-four} is an existence result for a fixed finite
description. Its proof does not supply an efficient multiplier-selection
algorithm or coefficient-size bounds. The exact checks accompanying this
section verify the finite identities and witnesses in
Example~\ref{ex:four-sharp}; the perturbation and limiting arguments are
proved above rather than inferred from those checks.
